\section{Models considered}
\label{sec:formatting}

All text must be in a two-column format.
The total allowable size of the text area is $6\frac78$ inches (17.46 cm) wide by $8\frac78$ inches (22.54 cm) high.
Columns are to be $3\frac14$ inches (8.25 cm) wide, with a $\frac{5}{16}$ inch (0.8 cm) space between them.
The main title (on the first page) should begin 1 inch (2.54 cm) from the top edge of the page.
The second and following pages should begin 1 inch (2.54 cm) from the top edge.
On all pages, the bottom margin should be $1\frac{1}{8}$ inches (2.86 cm) from the bottom edge of the page for $8.5 \times 11$-inch paper;
for A4 paper, approximately $1\frac{5}{8}$ inches (4.13 cm) from the bottom edge of the
page.

%-------------------------------------------------------------------------
\subsection{EfficientNet}
literature: arXiv:1905.11946 [cs.LG] EfficientNet: Rethinking Model Scaling for Convolutional Neural Networks

EfficientNet is a convolutional neural network (CNN) architecture that introduces a systematic scaling method to optimize model performance. 
Instead of arbitrarily increasing network depth, width, or input resolution, EfficientNet uniformly scales these three dimensions using a compound scaling coefficient. 
This approach is based on the following intuitions:
\begin{itemize}
\item {\bf Larger input images require greater depth} to capture long-range dependencies 
\item {\bf Increased width} is needed to learn more detailed and fine-grained features 
\item {\bf Higher resolution} inputs should be accompanied by proportional increases in both depth and width in order to maintain a balanced and efficient network structure
\end{itemize}

By coordinating these scaling dimensions through fixed scaling factors, EfficientNet achieves a better trade-off between accuracy and efficiency compared to traditional CNN architectures.
For the competition, we performed feature extraction using EfficientNetB0 and EfficientNetB4 models, both pretrained on ImageNet and accessed via TensorFlow's keras.applications module. 
To reduce the spatial dimensions of the output feature maps, we applied a Global Average Pooling (GAP) layer, which averages the value of each feature map, transforming a 3D tensor into a 1D vector of fixed length that is more suitable for similarity comparison. 
For the retrieval task, we used Annoy (Approximate Nearest Neighbors Oh Yeah) to perform fast and memory-efficient nearest neighbor searches. Although the initial implementations used Euclidean distance, we transitioned to cosine similarity, which better suits the nature of deep image embeddings. This was implemented by L2-normalizing the extracted feature vectors before indexing them, ensuring that angular distance (as used by Annoy) corresponds to cosine distance.

Annoy is particularly well-suited for this use case because:
\begin{itemize}
\item it supports {\bf angular distance} which is well suited for normalized embeddings 
\item it is {\bf highly efficient and scalable}, even with high dimensional data 
\item enables {\bf fast retrieval times} with low memory usage thanks to the tree-based structure. 
\end{itemize}

The attempt to fine-tune EfficientNet did not provide meaningful improvement in model's performance, which is why it was not tested on competition day. 

%-------------------------------------------------------------------------
\subsection{ResNet}

ResNet (Residual Networks) is a deep convolutional neural network architecture and it introduces residual (or skip) connections, which allow gradients to flow directly thorugh the network, enabling training of much deeper architecture.
OOur investigation employed three distinct ResNet-based approaches for content-based image retrieval, each implementing different transfer learning paradigms and similarity computation methodologies.
We utilized a pre-trained ResNet50 model as a fixed feature extractor, leveraging representations learned from ImageNet classification. 
The network's final fully connected layer was replaced with an identity mapping, transforming the model from a classifier to a feature encoder. This modification enabled extraction of 2048-dimensional feature vectors from the global average pooling layer preceding the original classification head.
L2-normalized feature vectors were compared using cosine distance to find the 10 most similar images to the query. 
The second methodology employed ResNet18 architecture with full network fine-tuning using cross-entropy loss for multi-class classification on the target domain. Following training convergence, the classification head was removed, and 512-dimensional feature vectors were extracted from the penultimate layer.
Similarity search was implemented using Facebook's FAISS (Facebook AI Similarity Search) library with L2 (Euclidean) distance indexing.
The third strategy implemented selective fine-tuning of ResNet50, where all convolutional layers remained frozen (requires _ grad=False) while only the final classification layer underwent parameter updates. 
This conservative transfer learning approach preserves the robust feature representations learned from ImageNet while adapting the decision boundary to the target classification task.
After training, the model was suited for feature extraction (by removing the trained classification layer), leading to 2048-dimensional normalized vectors used for similarity computations using cosine distance. 

%-------------------------------------------------------------------------
\subsection{VGG}

VGG-16 is a classical convolutional architecture known for its simplicity and uniform design, consisting of sequential 3×3 convolutional layers followed by max-pooling and fully connected layers~\cite{simonyan2015very}.
We decided to try it because of its layer-wise simplicity and consistent use of small convolutional filters make it interpretable and stable across a variety of visual tasks.
We included VGG-16 as a diagnostic baseline to validate our retrieval pipeline and to assess how well a classical CNN performs on a modern cross-domain retrieval task.
However, in our retrieval scenario, it exhibited several limitations. Its strong reliance on local texture patterns made it susceptible to the types of superficial variations introduced by synthetic image generation — such as lighting inconsistencies, skin smoothness artifacts, or uniform backgrounds.
Additionally, the absence of skip connections limited its ability to integrate deep semantic context, which proved essential for cross-domain identity recognition.
Retrieval rankings often favored images with similar lighting or pose over those with true identity matches. These results aligned with prior findings that CNNs trained on ImageNet are overly texture-biased~\cite{geirhos2019imagenet},

%-------------------------------------------------------------------------
\subsection{DINO V2}

Meta's DINOv2~\cite{caron2021emerging} is a recent self-supervised vision transformer (ViT) trained using a teacher-student framework that avoids the need for manual labels. Unlike traditional supervised learning, DINO trains models to maximize consistency across different augmented views of the same image.
This promotes the learning of features that are robust to low-level variations such as texture, lighting, and background, and encourages a focus on object-level or identity-preserving semantics.
For those reasons it showcases an ability to generalize across different visual domains without the need for fine tuning, which made it an obvious choice, given the task.
Specifically we used the vit_base_patch16_224.dino model. This corresponds to a Vision Transformer (ViT-B/16) trained using the original DINO self-supervised framework~\cite{caron2021emerging}, due to its capacity to form semantically coherent representations without requiring human-labeled supervision.
Neverthless, we obtained very low performances during the competition, probably because Dino is not trained to distinguish between identities and it lacks explicit supervision for facial or fine-grained matching tasks, as argued by Oquab et al.~\cite{oquab2023dinov2}.
Additionally, adding a fine-tuning layer did not help; in fact, the performance decreased. Our most likely explanation for such a result is that being pre-trained with unlabeled semantic relationships, fine-tuning it on a labeled training set might distort its geometry~\cite{oquab2023dinov2}.


%-------------------------------------------------------------------------
\subsection{CLIP}

The CLIP (Contrastive Language–Image Pretraining) framework, developed by OpenAI, was the most effective model we evaluated. CLIP is trained on hundreds of millions of image-text pairs scraped from the internet, learning to align images and natural language descriptions in a shared embedding space.
This multimodal contrastive training equips CLIP with a robust understanding of high-level semantics, making it especially well-suited for retrieval tasks where visual similarity is more conceptual~\cite{radford2021learning}.
We tested several variants of CLIP in frozen mode, including ViT-B/16, ViT-L/14, and RN50x64. All models were evaluated by computing cosine similarity between image embeddings. Among them, RN50x64 obtained the best performance during the competition window.
Another factor in the choice to deploy it was that it is trained on a wide range of image-text pairs during training, CLIP generalizes well to new domains and unseen classes and it should also achieve high performances in unlabeled settings.
As seen, CLIP models come in multiple architectural forms — most notably Vision Transformers (ViTs) and ResNet-based (RN) backbones. We chose to evaluate both because each brings different inductive biases and practical trade-offs that can influence retrieval performance depending on the task.
While ViT based models are particularly adept at capturing long-range semantic relationships and context-aware representations~\cite{dosovitskiy2021an}, ResNet based models are often more robust to small-scale distortions, faster to train or fine-tune, and more efficient in memory-constrained environments. Their convolutional structure also tends to preserve fine-grained spatial details that are useful in identity-level retrieval~\cite{he2016deep}.
We decided to test both because of the unknown nature of the retrieval task we would face.
During competition, we choose to fine-tune only RN50x64 for time issues, and we choose it because we hypotized that ResNets would be more stable during fine-tuning, especially on small datasets; also RN50x64 is significantly more memory-efficient than ViT-L/14, which was key given the context.


\begin{table}
  \centering
  \begin{tabular}{@{}lc@{}}
    \toprule
    Method & Frobnability \\
    \midrule
    Theirs & Frumpy \\
    Yours & Frobbly \\
    Ours & Makes one's heart Frob\\
    \bottomrule
  \end{tabular}
  \caption{Results.   Ours is better.}
  \label{tab:example}
\end{table}

%-------------------------------------------------------------------------
\subsection{Illustrations, graphs, and photographs}

All graphics should be centered.
In \LaTeX, avoid using the \texttt{center} environment for this purpose, as this adds potentially unwanted whitespace.
Instead use
{\small\begin{verbatim}
  \centering
\end{verbatim}}
at the beginning of your figure.
Please ensure that any point you wish to make is resolvable in a printed copy of the paper.
Resize fonts in figures to match the font in the body text, and choose line widths that render effectively in print.
Readers (and reviewers), even of an electronic copy, may choose to print your paper in order to read it.
You cannot insist that they do otherwise, and therefore must not assume that they can zoom in to see tiny details on a graphic.

When placing figures in \LaTeX, it's almost always best to use \verb+\includegraphics+, and to specify the figure width as a multiple of the line width as in the example below
{\small\begin{verbatim}
   \usepackage{graphicx} ...
   \includegraphics[width=0.8\linewidth]
                   {myfile.pdf}
\end{verbatim}
}


%-------------------------------------------------------------------------
\subsection{Color}

Please refer to the author guidelines on the \confName\ \confYear\ web page for a discussion of the use of color in your document.

If you use color in your plots, please keep in mind that a significant subset of reviewers and readers may have a color vision deficiency; red-green blindness is the most frequent kind.
Hence avoid relying only on color as the discriminative feature in plots (such as red \vs green lines), but add a second discriminative feature to ease disambiguation.